%! Introducing Statistics: What Can We Learn from Data? — Unit 1, Lesson 1.0 — Slides (Beamer)
\documentclass[aspectratio=169, 11pt]{beamer}
\usepackage{apstats-beamer}   % provides \navyheader{} and \sectionlabel[color]{}
\usetikzlibrary{positioning}  % -beamer loads tikz but not this library

% The C and Y column types live in apstats-article, which a beamer deck does not
% load — redeclare them here rather than pulling in the article preamble.
\newcolumntype{C}[1]{>{\centering\arraybackslash}p{#1}}
\newcolumntype{Y}{>{\centering\arraybackslash}X}

\begin{document}

% TITLE SLIDE (CourseName is not defined in beamer, so the name is literal)
{
\setbeamercolor{background canvas}{bg=navy}
\begin{frame}[plain]
  \vspace{2.8cm}\hspace{0.55cm}%
  \begin{minipage}{0.9\textwidth}
    {\color{goldacc}\bfseries\small LESSON 1.0}\\[0.5cm]
    {\color{white}\bfseries\LARGE Introducing Statistics: What Can We Learn from Data?}\\[1.2cm]
    {\color{skymid}\small Statistics~$\cdot$~Unit 1}
  \end{minipage}
\end{frame}
}

% The deck follows the EFFL flow: targets (SPOILER-FREE) → warm-up → activity launch →
% debrief frames (formalize in "red ink") → QuickNotes arc → application → practice →
% homework. Everything before the debrief divider stays vocabulary-free.
\newcommand{\redink}[1]{{\color{redacc}\bfseries #1}}

% ── LEARNING TARGETS (spoiler-free: no `data', `variable', `statistical question') ──
\begin{frame}[t, plain]
  \navyheader{Today: experience first, formalize later}
  \sectionlabel{Learning targets}
  By the end of today, I can\ldots
  \begin{itemize}\setlength{\itemsep}{5pt}
    \item explain why a number on its own tells you nothing.
    \item read a table and say what one row stands for and what each column keeps track of.
    \item tell apart questions you look up and questions that need a whole group.
    \item write a question of my own that recorded information could answer.
  \end{itemize}
  \begin{block}{How today works}
    Warm-up $\to$ \textbf{the adoption board} with your group $\to$ we name what you found
    $\to$ practice. You will figure it out \emph{before} I give you the words.
  \end{block}
\end{frame}

% ── HOOK ─────────────────────────────────────────────────────────────────────
{
\setbeamercolor{background canvas}{bg=navy}
\begin{frame}[plain]
  \vspace{1.6cm}
  \begin{center}
    {\color{skymid}\bfseries\small WHAT DOES THIS TELL YOU?}\\[1.0cm]
    {\color{white}\bfseries\fontsize{72}{72}\selectfont 68}\\[1.2cm]
    {\color{goldacc}\small That is the whole slide. That is all you get.}
  \end{center}
\end{frame}
}

% ── ACTIVITY LAUNCH ──────────────────────────────────────────────────────────
\begin{frame}[t, plain]
  \navyheader{Experience \& Formalize: The Adoption Board}
  \sectionlabel[goldacc]{Groups of 3--4 $\cdot$ 20 minutes}
  One shelter board. Visitors read it; so does the director --- not for the same things.
  \begin{itemize}\setlength{\itemsep}{5pt}
    \item \textbf{Problem 1} --- read the board.
    \item \textbf{Problem 2} --- four questions the director might ask.
  \end{itemize}
  \begin{block}{Your job}
    Use only what you already know. For \emph{every} answer, be ready to show me
    \textbf{where on the board} you see it.
  \end{block}
\end{frame}

% ── DEBRIEF DIVIDER — everything after this point may name the vocabulary ─────
{
\setbeamercolor{background canvas}{bg=navy}
\begin{frame}[plain]
  \vspace{1.8cm}\hspace{0.8cm}%
  \begin{minipage}{0.86\textwidth}
    {\color{goldacc}\bfseries\small NOW WE NAME IT}\\[0.7cm]
    {\color{white}\bfseries\Large Everything on the next slides, you already found.
    We are only giving it the words statisticians use.}\\[0.9cm]
    {\color{skymid}\small Fill in your QuickNotes box as each name lands.}
  \end{minipage}
\end{frame}
}

\begin{frame}[t, plain]
  \navyheader{A number with no context}
  \sectionlabel{it could be any of these}
  \begin{columns}[t]
    \begin{column}{0.48\textwidth}
      \begin{itemize}\setlength{\itemsep}{7pt}
        \item A test score, out of 100
        \item A temperature, in $^\circ$F
        \item A resting heart rate
        \item A height, in inches
        \item A speed, in miles per hour
      \end{itemize}
    \end{column}
    \begin{column}{0.48\textwidth}
      \begin{block}{Essential Knowledge (VAR-1.A.1)}
        Numbers may convey meaningful information \textbf{when placed in context}.
      \end{block}
      \vspace{4pt}
      {\small A number is not data until you can say \textbf{who} it describes,
      \textbf{what} was measured, and \textbf{in what units}.}
    \end{column}
  \end{columns}
\end{frame}

% ── DEFINITIONS ──────────────────────────────────────────────────────────────
\begin{frame}[t, plain]
  \navyheader{Statistics}
  \sectionlabel{the definition we will use all year}
  \begin{block}{}
    \textbf{Statistics} is the science of \textbf{collecting}, \textbf{organizing},
    \textbf{analyzing}, and \textbf{drawing conclusions from} data.
  \end{block}
  \vspace{10pt}
  \begin{itemize}\setlength{\itemsep}{8pt}
    \item Note what is \emph{not} in that definition: arithmetic.
    \item The computation is the smallest part. The question and the context
          are the work.
    \item \textbf{Data} are recorded values \emph{together with} the context that
          gives them meaning.
  \end{itemize}
\end{frame}

\begin{frame}[t, plain]
  \navyheader{The W questions}
  \sectionlabel{answer these and a number becomes data}
  \vspace{4pt}
  \renewcommand{\arraystretch}{1.6}
  \begin{tabularx}{\textwidth}{@{} >{\bfseries\color{navy}}l @{\hspace{1.2cm}} X @{}}
    \hline
    Who?           & Seven students in this class \\
    What?          & Their height \\
    Units?         & Inches \\
    When \& Where? & Today, in this classroom \\
    \hline
  \end{tabularx}
  \vspace{12pt}
  \begin{center}
    {\color{navy}\bfseries 68 \quad 71 \quad 64 \quad 70 \quad 66 \quad 73 \quad 69}\\[6pt]
    {\small The same seven numbers from the warm-up --- now they answer something.}
  \end{center}
\end{frame}

% ── INDIVIDUALS AND VARIABLES ────────────────────────────────────────────────
\begin{frame}[t, plain]
  \navyheader{Individuals and variables}
  \vspace{2pt}
  \renewcommand{\arraystretch}{1.3}
  \begin{tabularx}{\textwidth}{@{} l Y Y Y Y @{}}
    \rowcolor{sky}
    \textbf{Name} & \textbf{Breed} & \textbf{Age (yr)} & \textbf{Weight (lb)} & \textbf{Adopted?} \\
    \hline
    Biscuit & Beagle    & 3 & 24 & Yes \\
    Momo    & Terrier   & 1 & 15 & No  \\
    Ranger  & Shepherd  & 5 & 61 & Yes \\
    Pepper  & Retriever & 2 & 55 & No  \\
    \hline
  \end{tabularx}
  \vspace{12pt}
  \begin{columns}[t]
    \begin{column}{0.48\textwidth}
      \textbf{\color{navy}Each row} is one \textbf{individual}.\\[4pt]
      {\small The person, animal, or object described.}
    \end{column}
    \begin{column}{0.48\textwidth}
      \textbf{\color{navy}Each column} is one \textbf{variable}.\\[4pt]
      {\small A characteristic recorded for every individual.}
    \end{column}
  \end{columns}
\end{frame}

\begin{frame}[t, plain]
  \navyheader{A variable is not a value}
  \sectionlabel{the most common day-one error}
  \vspace{6pt}
  \begin{columns}[t]
    \begin{column}{0.46\textwidth}
      \begin{block}{Variable}
        \textbf{Breed} --- the column heading. The characteristic being recorded.
      \end{block}
    \end{column}
    \begin{column}{0.46\textwidth}
      \begin{block}{Value}
        \textbf{Beagle} --- what one individual happens to be. An entry in the column.
      \end{block}
    \end{column}
  \end{columns}
  \vspace{14pt}
  \begin{center}
    {\large ``Beagle'' is a \textbf{value} of the variable \textbf{Breed}.}\\[8pt]
    {\small\color{slate} Point at the column heading before you answer.}
  \end{center}
\end{frame}

% ── VARIABILITY ──────────────────────────────────────────────────────────────
\begin{frame}[t, plain]
  \navyheader{Why does statistics exist?}
  \sectionlabel{the reason for the entire course}
  \vspace{6pt}
  \begin{itemize}\setlength{\itemsep}{10pt}
    \item If every dog weighed exactly the same, one dog would tell you everything
          about weight.
    \item You would not need a sample, a graph, or an average. You would need
          one measurement.
    \item But the values \textbf{differ} --- 24, 15, 61, 55 --- so one measurement
          is not enough.
  \end{itemize}
  \vspace{10pt}
  \begin{block}{Variability}
    Values of a variable \textbf{differ from one individual to the next}. Variability
    is the reason we collect data at all.
  \end{block}
\end{frame}

% ── STATISTICAL QUESTIONS ────────────────────────────────────────────────────
\begin{frame}[t, plain]
  \navyheader{What makes a question statistical?}
  \sectionlabel{both checks must pass}
  \vspace{8pt}
  \begin{columns}[t]
    \begin{column}{0.46\textwidth}
      \begin{block}{Check 1}
        Will the answers \textbf{differ} across individuals?
      \end{block}
    \end{column}
    \begin{column}{0.46\textwidth}
      \begin{block}{Check 2}
        Would I have to \textbf{collect data} to answer it?
      \end{block}
    \end{column}
  \end{columns}
  \vspace{16pt}
  \begin{center}
    {\large Both yes $\Rightarrow$ \textbf{\color{navy}statistical question.}}\\[8pt]
    {\small One definite answer $\Rightarrow$ not statistical.}
  \end{center}
\end{frame}

\begin{frame}[t, plain]
  \navyheader{S or N?}
  \sectionlabel{call it out}
  \vspace{4pt}
  \renewcommand{\arraystretch}{1.7}
  \begin{tabularx}{\textwidth}{@{} X C{3.2cm} @{}}
    \hline
    \textbf{Question} & \textbf{S or N} \\
    \hline
    How many minutes did I sleep last night?          & \\
    How many minutes do students at our school sleep? & \\
    What is the weight of Ranger, the shepherd?       & \\
    How do the weights of shelter dogs vary?          & \\
    \hline
  \end{tabularx}
  \vspace{10pt}
  \begin{center}
    {\small\color{slate} For every \textbf{N}: what one change would make it statistical?}
  \end{center}
\end{frame}

% ── THE ARC ──────────────────────────────────────────────────────────────────
\begin{frame}[t, plain]
  \navyheader{The data analysis process}
  \sectionlabel{this is the shape of the whole course}
  \vspace{18pt}
  \begin{center}
  \begin{tikzpicture}[node distance=0.5cm, every node/.style={font=\small\bfseries}]
    \node[fill=navy, text=white, rounded corners=2mm, minimum height=1.1cm,
          minimum width=2.6cm] (a) {Ask a question};
    \node[fill=skymid, text=charcoal, rounded corners=2mm, minimum height=1.1cm,
          minimum width=2.6cm, right=of a] (b) {Collect data};
    \node[fill=goldacc, text=charcoal, rounded corners=2mm, minimum height=1.1cm,
          minimum width=2.6cm, right=of b] (c) {Analyze};
    \node[fill=greenacc, text=white, rounded corners=2mm, minimum height=1.1cm,
          minimum width=2.6cm, right=of c] (d) {Interpret};
    \draw[->, thick, navy] (a) -- (b);
    \draw[->, thick, navy] (b) -- (c);
    \draw[->, thick, navy] (c) -- (d);
  \end{tikzpicture}
  \end{center}
  \vspace{16pt}
  \begin{center}
    {\small Today is step one. Every answer you write this year ends back in
    \textbf{context}.}
  \end{center}
\end{frame}

% ── APPLICATION ──────────────────────────────────────────────────────────────
\begin{frame}[t, plain]
  \navyheader{Application: The Gym Sign-Up}
  \sectionlabel[goldacc]{We work this one together --- you hold the pen}
  A gym records the \textbf{resting heart rate} (bpm) of each of its 120 members on the
  morning they join, and each member's \textbf{home ZIP code}.
  \begin{itemize}\setlength{\itemsep}{5pt}
    \item Who is one individual here --- the gym, or a member?
    \item Write a statistical question, and say what \emph{varies} in it.
    \item ``What is the most common ZIP code?'' --- statistical or not?
  \end{itemize}
  \begin{block}{Run both checks before you answer}
    Does check 1 say the answers have to be \emph{numbers}, or just that they have to
    \redink{differ}?
  \end{block}
\end{frame}

% ── PRACTICE (in class, not scored) ──────────────────────────────────────────
\begin{frame}[t, plain]
  \navyheader{Check Your Understanding}
  \sectionlabel{5 exercises $\cdot$ pairs first, then on your own}
  \begin{itemize}\setlength{\itemsep}{5pt}
    \item Turtles in a pond: individuals, variable, units.
    \item Write your own statistical question --- and name what varies.
    \item Same pond, two questions: which one is statistical, and why?
    \item ``The average is 12.'' What are you missing?
    \item One multiple-choice item, with a one-sentence justification.
  \end{itemize}
  \begin{block}{This is practice, not a quiz}
    These are \textbf{not scored}. Start with 1, 3, and 5. I will look over your shoulder
    at 5 on the way out.
  \end{block}
\end{frame}

% ── HOMEWORK ─────────────────────────────────────────────────────────────────
\begin{frame}[t, plain]
  \navyheader{AP Practice, homework \& what's next}
  \sectionlabel[goldacc]{AP Practice --- practice, not a grade}
  \begin{itemize}\setlength{\itemsep}{5pt}
    \item Four multiple-choice items on the adoption board and a vet clinic.
    \item One free-response set: a school-newspaper survey of the senior class.
    \item Optional: bring in a headline with a number and strip it for context.
  \end{itemize}
  \begin{block}{Your homework is in DeltaMath}
    The \redink{graded} assignment is posted in \textbf{DeltaMath} --- log in and complete it
    before next class. The AP Practice page is not scored: work it, and bring what you get
    stuck on to study hall or office hours.
  \end{block}
\end{frame}

% ── CLOSE ────────────────────────────────────────────────────────────────────
{
\setbeamercolor{background canvas}{bg=navy}
\begin{frame}[plain]
  \vspace{1.5cm}\hspace{0.8cm}%
  \begin{minipage}{0.86\textwidth}
    {\color{goldacc}\bfseries\small BEFORE YOU GO}\\[0.7cm]
    {\color{white}\bfseries\Large A number is not data until you can say who it
    describes, what was measured, and in what units.}\\[0.9cm]
    {\color{skymid}\small Next: Lesson 1.1 --- categorical vs.\ quantitative variables,
    the first decision that shapes every graph we draw.}
  \end{minipage}
\end{frame}
}

\end{document}
